\documentclass[11pt]{article}
\usepackage[a4paper,margin=1in]{geometry}
\usepackage{fontspec}
\usepackage[boldfont=false]{xeCJK}
\setCJKmainfont{FandolSong-Regular.otf}
\usepackage{amsmath}
\usepackage{amssymb}
\usepackage{array}
\usepackage{booktabs}
\usepackage{graphicx}
\usepackage{adjustbox}
\usepackage{enumitem}
\setlist[itemize]{topsep=2pt,partopsep=0pt,parsep=0pt,itemsep=2pt,leftmargin=1.4em}
\setlist[enumerate]{topsep=2pt,partopsep=0pt,parsep=0pt,itemsep=2pt,leftmargin=1.6em}
\usepackage{parskip}
\usepackage{fvextra}
\fvset{breaklines=true,breakanywhere=true,fontsize=\small}
\setlength{\parindent}{0pt}

\begin{document}

\begin{center}
  {\LARGE\bfseries The price system and the microeconomy (A Level)}\\[0.35em]
  {\large A Level Economics}
\end{center}
\vspace{0.6em}

This topic looks more deeply at consumers, costs and firms.

\section*{Utility and consumer choice}

\textbf{utility} 效用 is the satisfaction a consumer gets from a good.

\begin{itemize}
\item \textbf{total utility} 总效用 is the satisfaction from all the units consumed.

\item \textbf{marginal utility} 边际效用 is the extra satisfaction from one more unit.

\end{itemize}

The law of \textbf{diminishing marginal utility} 边际效用递减 says each extra unit gives less extra satisfaction. Total utility still rises, but more and more slowly.

\subsection*{The equi-marginal principle}

A consumer with a fixed budget gets the most total utility when the last dollar spent on each good gives the same extra utility. This is the \textbf{equi-marginal principle} 等边际原则:

\[
\frac{MU_A}{P_A} = \frac{MU_B}{P_B}
\]

Because marginal utility falls as you buy more, you only buy extra units when the price falls. This is why the \textbf{demand curve} 需求曲线 slopes downward.

\subsection*{Limitations}

The approach is hard to use in practice: utility cannot really be measured in numbers, and people do not calculate like this. It also ignores habit and advertising.

\section*{Indifference curves and budget lines}

An \textbf{indifference curve} 无差异曲线 shows all the combinations of two goods that give a consumer the \textbf{same} total utility. It slopes downward (more of one good means less of the other) and bows in towards the origin.

A \textbf{budget line} 预算线 shows all the combinations the consumer can just afford with their income, at the current prices.

\textbf{consumer equilibrium} 消费者均衡 is the best the consumer can do: the point where the budget line just touches the highest possible indifference curve.

\subsection*{Income and substitution effects}

When the price of a good falls, the consumer buys more for two reasons:

\begin{itemize}
\item the \textbf{substitution effect} 替代效应 — the good is now cheaper than others, so the consumer switches towards it.

\item the \textbf{income effect} 收入效应 — the consumer's money now buys more, so real income has risen.

\end{itemize}

These effects work differently for different goods:

\begin{itemize}
\item for a \textbf{normal good} 正常品, both effects raise the quantity bought.

\item for an \textbf{inferior good} 低档品, the income effect works the other way (higher real income lowers demand), but it is usually weaker, so demand still rises when price falls.

\item a \textbf{Giffen good} 吉芬品 is a rare, very inferior good where the income effect is so strong that demand actually \textbf{falls} when the price falls. Its demand curve slopes upward.

\end{itemize}

\section*{Efficiency and market failure}

Economists judge how well resources are used.

\begin{itemize}
\item \textbf{productive efficiency} 生产效率 — making goods at the lowest possible cost (using the fewest resources).

\item \textbf{allocative efficiency} 配置效率 — making the goods people most want, where price equals marginal cost.

\item \textbf{Pareto optimality} 帕累托最优 — a point where no one can be made better off without making someone else worse off.

\item \textbf{dynamic efficiency} 动态效率 — improving over time through new technology and better products.

\end{itemize}

\textbf{market failure} 市场失灵 happens when the free market does not reach an efficient result. The main causes are monopoly power, \textbf{externalities} 外部性, \textbf{public goods} 公共物品, \textbf{information failure} 信息失灵, and \textbf{factor immobility} 要素不流动性 (factors cannot move quickly to where they are needed).

\section*{Externalities and social costs and benefits}

Some costs and benefits fall on people outside the deal. We separate three levels:

\begin{itemize}
\item a \textbf{private cost} 私人成本 or \textbf{private benefit} 私人收益 falls on the buyer or seller.

\item an \textbf{external cost} 外部成本 or \textbf{external benefit} 外部收益 falls on a third party.

\item the \textbf{social cost} 社会成本 is private cost plus external cost; the \textbf{social benefit} 社会收益 is private benefit plus external benefit.

\end{itemize}

\subsection*{Marginal analysis}

We study externalities with \textbf{marginal analysis} 边际分析, using these curves:

\begin{itemize}
\item \textbf{marginal private cost} 边际私人成本 (MPC) and \textbf{marginal social cost} 边际社会成本 (MSC).

\item \textbf{marginal private benefit} 边际私人收益 (MPB) and \textbf{marginal social benefit} 边际社会收益 (MSB).

\end{itemize}

The best outcome for society is where MSC = MSB. But the market produces where MPC = MPB. The gap between these creates lost welfare.

\begin{itemize}
\item a \textbf{negative externality} 负外部性 (like pollution) means MSC is above MPC. The market makes \textbf{too much}.

\item a \textbf{positive externality} 正外部性 (like education) means MSB is above MPB. The market makes \textbf{too little}.

\end{itemize}

The lost welfare from making too much or too little is the \textbf{deadweight welfare loss} 无谓损失 — the triangle on the diagram between the social and private curves.

\section*{Costs, revenue and profit}

\subsection*{Costs}

In the \textbf{short run} 短期, at least one factor is fixed; in the \textbf{long run} 长期, all factors can change.

\begin{center}
\adjustbox{max width=\linewidth}{%
\begin{tabular}{ll}
\toprule
\textbf{Cost} & \textbf{Meaning} \\
\midrule
\textbf{fixed cost} 固定成本 & does not change with output (rent on a factory) \\
\textbf{variable cost} 可变成本 & changes with output (raw materials) \\
\textbf{total cost} 总成本 & fixed cost plus variable cost \\
\textbf{average cost} 平均成本 & total cost per unit: $AC = TC / Q$ \\
\textbf{marginal cost} 边际成本 & the cost of one more unit: $MC = \Delta TC / \Delta Q$ \\
\bottomrule
\end{tabular}}
\end{center}

In the short run, output is limited by the law of \textbf{diminishing returns} 边际报酬递减: as you add more of a variable factor (workers) to a fixed factor (machines), each extra worker eventually adds less output. This makes marginal cost rise.

In the long run, all factors vary, so we look at \textbf{returns to scale} 规模报酬. As a firm grows, it can gain \textbf{economies of scale} 规模经济 (lower average cost from buying in bulk, using big machines, and specialising). If it grows too big, it may suffer \textbf{diseconomies of scale} 规模不经济 (poor communication and control), and average cost rises.

\subsection*{Revenue}

\begin{center}
\adjustbox{max width=\linewidth}{%
\begin{tabular}{ll}
\toprule
\textbf{Revenue} & \textbf{Meaning} \\
\midrule
\textbf{total revenue} 总收益 & all money from sales: $TR = P \times Q$ \\
\textbf{average revenue} 平均收益 & revenue per unit, which equals the price \\
\textbf{marginal revenue} 边际收益 & the revenue from one more unit: $MR = \Delta TR / \Delta Q$ \\
\bottomrule
\end{tabular}}
\end{center}

\subsection*{Profit}

Profit is total revenue minus total cost. Economists include the opportunity cost of the owner's resources in cost, so:

\begin{itemize}
\item \textbf{normal profit} 正常利润 is just enough profit to keep the firm in the industry. It counts as a cost.

\item \textbf{supernormal profit} 超常利润 (also called abnormal profit) is profit above normal.

\item \textbf{subnormal profit} 次正常利润 is profit below normal — the firm may leave in the long run.

\end{itemize}

A firm makes the most profit at the output where $MR = MC$.

\section*{Market structures}

A \textbf{market structure} 市场结构 describes a market by the number of firms, the barriers to entry, and the type of product.

\begin{center}
\adjustbox{max width=\linewidth}{%
\begin{tabular}{lllll}
\toprule
\textbf{Structure} & \textbf{Firms} & \textbf{Barriers} & \textbf{Product} & \textbf{Long-run profit} \\
\midrule
\textbf{perfect competition} 完全竞争 & very many & none & identical & normal only \\
\textbf{monopolistic competition} 垄断竞争 & many & low & slightly different & normal \\
\textbf{oligopoly} 寡头垄断 & a few & high & varies & often supernormal \\
\textbf{monopoly} 垄断 & one & very high & unique & supernormal \\
\bottomrule
\end{tabular}}
\end{center}

\begin{itemize}
\item in \textbf{perfect competition}, each firm is a tiny \textbf{price taker} 价格接受者 — it must accept the market price. Firms make only normal profit in the long run, and the market is both productively and allocatively efficient.

\item in \textbf{monopolistic competition}, many firms sell slightly different products (\textbf{product differentiation} 产品差异化), so each has a little price-setting power, but free entry competes profit away.

\item in an \textbf{oligopoly}, a few large firms dominate. They watch each other closely. They may compete hard on price, or they may agree to fix prices — this is \textbf{collusion} 合谋, and a group that does so is a \textbf{cartel} 卡特尔.

\item a \textbf{monopoly} is a single seller protected by high \textbf{barriers to entry} 进入壁垒. It can earn supernormal profit for a long time but is usually inefficient.

\end{itemize}

\subsection*{Price and non-price competition}

\begin{itemize}
\item \textbf{price competition} 价格竞争 means cutting prices to win customers.

\item \textbf{non-price competition} 非价格竞争 means competing in other ways — advertising, quality, branding and service.

\end{itemize}

\subsection*{Price discrimination}

\textbf{price discrimination} 价格歧视 is charging different consumers different prices for the \textbf{same} good (for example, cheaper train tickets for students). It needs price-setting power, groups with different elasticities, and a way to stop resale between groups. It raises the firm's profit.

\section*{The growth and survival of firms}

Firms grow in two ways:

\begin{itemize}
\item \textbf{internal growth} 内部增长 (organic) — growing bigger by selling more.

\item \textbf{external growth} 外部增长 — joining with another firm through a \textbf{merger} 合并 or takeover. Types:

\begin{itemize}
\item \textbf{horizontal integration} 横向一体化 — joining a firm at the same stage of the same industry (two car makers).

\item \textbf{vertical integration} 纵向一体化 — joining a firm at a different stage (a car maker buys a tyre maker).

\item \textbf{conglomerate integration} 混合一体化 — joining a firm in a different industry, to spread risk.

\end{itemize}

\end{itemize}

\subsection*{Why some firms stay small}

Many firms stay small because the market is small (local services), the owner wants to keep control, or there are few economies of scale. They survive by offering personal service or filling a \textbf{niche} 利基市场.

Firms also face \textbf{barriers to entry} and \textbf{barriers to exit} 退出壁垒 (such as costly machines that cannot be sold). High exit barriers can trap a firm in a loss-making market.

\section*{The objectives of firms}

Not every firm simply chases the most profit.

\begin{itemize}
\item \textbf{profit maximisation} 利润最大化 — producing where MR = MC. The traditional aim.

\item \textbf{revenue maximisation} 收益最大化 — producing where MR = 0, perhaps to grow market share.

\item \textbf{sales maximisation} 销售最大化 — selling as much as possible while still making normal profit.

\item \textbf{satisficing} 满意化 — aiming for "good enough" results that keep owners, workers and customers all reasonably happy.

\end{itemize}

\subsection*{Divorce of ownership from control}

In large companies the owners (shareholders) are not the managers. This \textbf{divorce of ownership from control} 所有权与控制权分离 creates the \textbf{principal-agent problem} 委托代理问题: managers (the agents) may pursue their own aims — bigger salaries or an easier life — instead of the owners' aim of maximum profit.


\end{document}
